\documentclass{article}
\usepackage{amsmath,amssymb}

\usepackage[margin=2cm]{geometry}

\begin{document}

\section*{Unitary Implementation of a dyadic LCU}

Let
\[
P(A) = \sum_{k=1}^K c_k\, U_k
\]
be a linear combination of unitary operators $U_k$ with real coefficients $c_k$
satisfying
\[
\sum_{k=1}^K |c_k| = 1.
\]

\medskip

To implement this operator using a uniform superposition, we introduce a
discretization parameter $m \in \mathbb{N}$ and define integers
\[
n_k \in \mathbb{N}
\quad \text{such that} \quad
c_k \approx \sigma(c_k) \frac{n_k}{2^m},
\qquad
\sum_{k=1}^K n_k = 2^m.
\]
where $\sigma(c_k)$ designates the sign of $c_k$.


Thus, each coefficient $c_k$ is approximated by a rational number with denominator $2^m$,
and $n_k$ represents the discrete weight (multiplicity) of $U_k$.

\medskip

Using the integers $n_k$, we define intervals
\[
I_k = [l_k, h_k),
\]
where
\[
l_1 = 0, \qquad
l_k = \sum_{j=1}^{k-1} n_j, \qquad
h_k = \sum_{j=1}^{k} n_j.
\]

These intervals form a partition
\[
\bigcup_{k=1}^K I_k = \{0,1,\dots,2^m - 1\},
\qquad
I_k \cap I_j = \varnothing \ \text{for } k \ne j,
\]
and satisfy
\[
|I_k| = n_k.
\]

\medskip



Then we obtain the identity
\[
\boxed{
\sum_{k=1}^K \sigma(c_k)\frac{n_k}{2^m}\, U_k
\approx
\sum_{k=1}^K c_k\, U_k.
}
\]

\medskip

This shows that the linear combination is realized as a uniform average
over the index register, where each unitary $U_k$ appears with multiplicity $n_k$.


\section*{Hilbert space structure}

We consider a Hilbert space of the form
\[
\mathcal{H}
= \mathcal{H}_x \otimes \mathcal{H}_b \otimes \mathcal{H}_{\mathrm{sys}},
\]
where:
\begin{itemize}
\item $\mathcal{H}_x$ is the address register (of size $2^m$),
\item $\mathcal{H}_b$ is a single ancilla qubit,
\item $\mathcal{H}_{\mathrm{sys}}$ is the system on which the operators $U_k$ act.
\end{itemize}


\section*{Operator primitives}

We assume the following unitary primitives:

\begin{enumerate}
\item \textbf{Comparator unitary}
\[
C_k \in \mathcal{U}(\mathcal{H}_x \otimes \mathcal{H}_b),
\]
defined by
\[
C_k \colon | x,b \rangle \mapsto | x,\, b \oplus \chi_{I_k}(x) \rangle,
\]
where $\chi_{I_k}(x)$ is the indicator function of the interval $I_k$.
So
$$
C_k = \bigoplus\limits_{x = 0,\dots,2^m - 1} X^{\chi_{I_k}(x)}
$$

\[
\chi_{I_k}(x) =
\begin{cases}
1, & x \in I_k, \\
0, & x \notin I_k.
\end{cases}
\]

This operator is reversible and hence unitary, with
\[
C_k^{-1} = C_k^\dagger.
\]

\item \textbf{Controlled operator}
\[
 C(\sigma(c_k)U_k) \in \mathcal{U}(\mathcal{H}_b \otimes \mathcal{H}_{\mathrm{sys}}),
\]
which acts as
\[
\text{``apply } \sigma(c_k)U_k \text{ if the ancilla is } 1,\ \text{otherwise } I.
\]
\end{enumerate}

\section*{Lifting to the full space}

We extend these operators to the full space:
\[
\widetilde{C}_k = C_k \otimes I_{\mathrm{sys}},
\]
\[
\widetilde{U}_k = I_x \otimes C(\sigma(c_k)U_k)
\]

\section*{Construction of a single block}

For each $k$, define the unitary operator
\[
S_k
= \widetilde{C}_k^\dagger
\circ \widetilde{U}_k
\circ \widetilde{C}_k.
\]

This operator has the following interpretation:

\begin{itemize}
\item $\widetilde{C}_k$ computes whether $x \in I_k$ into the ancilla,
\item $\widetilde{U}_k$ conditionally applies $\sigma(c_k)U_k$ to the system qubits.
\item $\widetilde{C}_k^\dagger$ uncomputes the ancilla.
\end{itemize}

Since it is a composition of unitary operators, $S_k$ is unitary.

\section*{Composit SELECT operator}

The SELECT operator is defined by composition:
\[
\boxed{
\mathrm{SELECT}
= S_1 \circ S_2 \circ \cdots \circ S_K.
}
\]
\section*{PREPARE operator}
The PREPARE operator is a simple diffuser
\[
\mathrm{PREPARE}\, |0\rangle
= \frac{1}{\sqrt{2^m}} \sum_{x=0}^{2^m-1} |x\rangle
\]
Implemented using a tensorproduct of Hadamard gates.
\end{document}